% ---
% title: C++17 十六进制浮点字面量
% date: 2026-07-13
% author: Crystal-Sky
% tags: C++, C++17, 基础语法, 浮点数, 字面量
% summary: 介绍十六进制浮点字面量的组成、二进制指数、类型后缀、计算方法与常见错误。
% slug: cpp-hexadecimal-floating-literals
% course: C++
% course_slug: cpp
% chapter: 基础语法
% chapter_order: 1
% lesson_order: 5
% lesson_type: 基础语法
% ---

\documentclass[UTF8,12pt,a4paper]{ctexart}

\usepackage{geometry}
\usepackage{xcolor}
\usepackage{listings}
\usepackage{amsmath}
\usepackage{booktabs}
\usepackage{enumitem}
\usepackage{tcolorbox}
\usepackage{fancyhdr}

\geometry{left=2.4cm,right=2.4cm,top=2.5cm,bottom=2.5cm}

\pagestyle{fancy}
\fancyhf{}
\fancyhead[L]{C++17 十六进制浮点字面量}
\fancyhead[R]{教学讲义}
\fancyfoot[C]{\thepage}

\definecolor{codebg}{RGB}{247,248,250}
\definecolor{keywordcolor}{RGB}{0,92,197}
\definecolor{commentcolor}{RGB}{0,128,0}
\definecolor{stringcolor}{RGB}{163,21,21}

\lstset{
    language=C++,
    basicstyle=\ttfamily\small,
    keywordstyle=\color{keywordcolor}\bfseries,
    commentstyle=\color{commentcolor},
    stringstyle=\color{stringcolor},
    backgroundcolor=\color{codebg},
    frame=single,
    numbers=left,
    numberstyle=\tiny\color{gray},
    numbersep=8pt,
    tabsize=4,
    showstringspaces=false,
    breaklines=true,
    columns=fullflexible,
    keepspaces=true
}

\newtcolorbox{conceptbox}[1]{
    colback=blue!3,
    colframe=blue!60!black,
    title=#1,
    fonttitle=\bfseries
}

\newtcolorbox{warningbox}[1]{
    colback=red!3,
    colframe=red!65!black,
    title=#1,
    fonttitle=\bfseries
}

\title{\textbf{C++17 十六进制浮点字面量}}
\author{}
\date{}

\begin{document}

\maketitle

\section{回顾普通浮点字面量}

平时常见的浮点字面量使用十进制表示：

\begin{lstlisting}
double a = 3.14;
double b = 1.5e2;   // 1.5 × 10^2 = 150
double c = 2.5e-3;  // 2.5 × 10^-3 = 0.0025
\end{lstlisting}

其中：

\begin{itemize}
    \item \texttt{e} 或 \texttt{E} 表示十进制指数；
    \item 指数部分的底数是 10。
\end{itemize}

\section{十六进制浮点字面量}

C++17 支持直接书写十六进制浮点字面量。

基本形式为：

\begin{lstlisting}
0x十六进制有效数字p二进制指数
\end{lstlisting}

例如：

\begin{lstlisting}
double x = 0x1.8p1;
\end{lstlisting}

它表示：

\[
0x1.8 \times 2^1
\]

这里需要注意：

\begin{itemize}
    \item 前缀使用 \texttt{0x} 或 \texttt{0X}；
    \item 有效数字使用十六进制；
    \item 指数标记必须使用 \texttt{p} 或 \texttt{P}；
    \item 指数的底数是 2，而不是 10。
\end{itemize}

\begin{conceptbox}{最重要的规则}
\[
\texttt{0xA.BpC}
=
\text{十六进制数 }0xA.B
\times 2^C
\]
\end{conceptbox}

\section{为什么指数用 \texttt{p}}

十六进制数字中，字母 \texttt{e} 本身就是合法数字：

\begin{lstlisting}
0x1e
\end{lstlisting}

这里的 \texttt{e} 表示十六进制数字 14。

因此，十六进制浮点数不能再使用 \texttt{e} 表示指数，只能使用
\texttt{p} 或 \texttt{P}。

\section{计算示例}

\subsection{\texttt{0x1p0}}

\begin{lstlisting}
double x = 0x1p0;
\end{lstlisting}

表示：

\[
1 \times 2^0 = 1
\]

所以 \texttt{x} 的值为 1。

\subsection{\texttt{0x1p3}}

\begin{lstlisting}
double x = 0x1p3;
\end{lstlisting}

表示：

\[
1 \times 2^3 = 8
\]

\subsection{\texttt{0x1.8p1}}

十六进制小数：

\[
0x1.8 = 1 + \frac{8}{16} = 1.5
\]

因此：

\[
0x1.8p1 = 1.5 \times 2^1 = 3
\]

\begin{lstlisting}
double x = 0x1.8p1;  // 3.0
\end{lstlisting}

\subsection{\texttt{0x1.4p2}}

先计算有效数字：

\[
0x1.4 = 1 + \frac{4}{16} = 1.25
\]

再乘以二进制指数：

\[
1.25 \times 2^2 = 5
\]

\begin{lstlisting}
double x = 0x1.4p2;  // 5.0
\end{lstlisting}

\subsection{负指数}

\begin{lstlisting}
double x = 0x1p-2;
\end{lstlisting}

表示：

\[
1 \times 2^{-2} = \frac{1}{4} = 0.25
\]


\section{类型后缀}

十六进制浮点字面量也可以使用浮点类型后缀。

\begin{lstlisting}
auto a = 0x1.8p1;   // double
auto b = 0x1.8p1f;  // float
auto c = 0x1.8p1L;  // long double
\end{lstlisting}

规则与普通浮点字面量一致：

\begin{center}
\begin{tabular}{ll}
\toprule
后缀 & 类型 \\
\midrule
无后缀 & \texttt{double} \\
\texttt{f} 或 \texttt{F} & \texttt{float} \\
\texttt{l} 或 \texttt{L} & \texttt{long double} \\
\bottomrule
\end{tabular}
\end{center}

\section{与普通浮点写法对比}

\begin{center}
\begin{tabular}{lll}
\toprule
写法 & 含义 & 结果 \\
\midrule
\texttt{1.5e2} & \(1.5 \times 10^2\) & 150 \\
\texttt{0x1.8p1} & \(1.5 \times 2^1\) & 3 \\
\texttt{0x1p4} & \(1 \times 2^4\) & 16 \\
\texttt{0x1p-3} & \(1 \times 2^{-3}\) & 0.125 \\
\bottomrule
\end{tabular}
\end{center}

\section{为什么有这种写法}

计算机中的浮点数通常采用二进制表示。

十进制小数有时不能被二进制精确表示，例如：

\begin{lstlisting}
double x = 0.1;
\end{lstlisting}

而十六进制浮点字面量与二进制浮点表示关系更直接。

例如：

\begin{lstlisting}
double x = 0x1p-3;
\end{lstlisting}

它准确表示：

\[
2^{-3} = 0.125
\]

因此十六进制浮点字面量常用于：

\begin{itemize}
    \item 精确表示二进制浮点值；
    \item 浮点算法测试；
    \item 底层数值程序；
    \item 需要明确控制浮点常量的场景。
\end{itemize}

\section{常见错误}

\subsection{使用 \texttt{e} 表示指数}

错误：

\begin{lstlisting}
double x = 0x1.8e2;
\end{lstlisting}

在十六进制中，\texttt{e} 会被当作数字的一部分。

正确写法：

\begin{lstlisting}
double x = 0x1.8p2;
\end{lstlisting}

\subsection{省略指数部分}

十六进制浮点字面量必须包含 \texttt{p} 指数部分。

错误：

\begin{lstlisting}
double x = 0x1.8;
\end{lstlisting}

正确：

\begin{lstlisting}
double x = 0x1.8p0;
\end{lstlisting}

\subsection{把指数理解为十进制指数}

\begin{lstlisting}
double x = 0x1p3;
\end{lstlisting}

它表示：

\[
1 \times 2^3
\]

不是：

\[
1 \times 10^3
\]



\section{总结}

\begin{conceptbox}{板书式记忆}
\begin{lstlisting}
0x1.8p1
\end{lstlisting}

表示：

\[
0x1.8 \times 2^1
\]

记住三点：

\begin{enumerate}[label=\arabic*.]
    \item \texttt{0x} 表示十六进制；
    \item \texttt{p} 表示指数；
    \item 指数的底数是 2。
\end{enumerate}
\end{conceptbox}

\end{document}
